\section{Obscure Questions in MVC Format}

\begin{exercise}
Let $X=\mathbb{R}^2\setminus\{(0,0)\}$. What is $\pi_1(X)$?

\begin{enumerate}[label=(\Alph*)]
    \item $0$
    \item $\mathbb{Z}$
    \item $\mathbb{Z}/2\mathbb{Z}$
    \item $\mathbb{Z}^2$
    \item A free group on two generators
\end{enumerate}
\end{exercise}
\begin{solution}
(B). $\R^2 \setminus \{(0,0)\}$ is homeomorphic to $\S^1$, which has 
first fundamental group $\Z$.
\end{solution}

\begin{exercise}
Let $G$ be a group of order $56$. Which of the following are the possible numbers of Sylow $7$-subgroups?

\begin{enumerate}[label=(\Alph*)]
    \item $1$ only
    \item $7$ only
    \item $1$ or $7$
    \item $1$ or $8$
    \item $7$ or $8$
\end{enumerate}
\end{exercise}
\begin{solution}
(D). Use Sylow III.
\end{solution}

\begin{exercise}
What is $\pi_1(\S^1\vee\S^1)$?
\begin{enumerate}[label=(\Alph*)]
    \item $\mathbb{Z}$
    \item $\mathbb{Z}^2$
    \item $\mathbb{Z}/2\mathbb{Z}$
    \item The free group on two generators
    \item $0$
\end{enumerate}
\end{exercise}
\begin{solution}
(D). This really requires the Seifert Van Kampen, but just remember it.
\end{solution}

\begin{exercise}
How many conjugacy classes are there in $S_5$?

\begin{enumerate}[label=(\Alph*)]
    \item $5$
    \item $6$
    \item $7$
    \item $8$
    \item $120$
\end{enumerate}
\end{exercise}
\begin{solution}
(C). This is just partitions of $5$ elements. There is:
\[
e, (12), (123), (1234), (12345), (12)(34), (123)(45)
\]
\end{solution}

\begin{exercise}
How many idempotent elements are there in $\mathbb{Z}/12\mathbb{Z}$? That is, how many $x$ satisfy $x^2=x$?

\begin{enumerate}[label=(\Alph*)]
    \item $1$
    \item $2$
    \item $3$
    \item $4$
    \item $6$
\end{enumerate}
\end{exercise}
\begin{solution}
(D). $0,1,4,9$. More generally, this follows by the CRT. $\Z/3\Z$ and $\Z/4\Z$ have two 
idempotents each.
\end{solution}

\begin{exercise}
Let $X=\mathbb{RP}^2$. Then $\pi_1(X)$ is isomorphic to
\begin{enumerate}[label=(\Alph*)]
    \item $0$
    \item $\mathbb{Z}$
    \item $\mathbb{Z}/2\mathbb{Z}$
    \item $\mathbb{Z}^2$
    \item $S_2$
\end{enumerate}
\end{exercise}
\begin{solution}
(C). Just remember it.
\end{solution}

\begin{exercise}
Suppose $G$ has order $21$. Which statement must be true?
\begin{enumerate}[label=(\Alph*)]
    \item Every subgroup is normal.
    \item $G$ has a unique Sylow $7$-subgroup.
    \item $G$ has a unique Sylow $3$-subgroup.
    \item $G$ is abelian.
    \item $G$ is cyclic.
\end{enumerate}
\end{exercise}
\begin{solution}
(B). Sylow I says that there are $3$-subgroups and $7$-subgroups.
Sylow III for $p = 7$ says that there is a unique Sylow $7$-subgroup.
\end{solution}

\begin{exercise}
Which space has fundamental group $\mathbb{Z}^2$?
\begin{enumerate}[label=(\Alph*)]
    \item $S^2$
    \item $S^1$
    \item $S^1\times S^1$
    \item $S^1\vee S^1$
    \item $D^2$
\end{enumerate}
\end{exercise}
\begin{solution}
(C). Products of spaces and first fundamental group commute.
\end{solution}

\begin{exercise}
How many units does $\mathbb{Z}/20\mathbb{Z}$ have?
\begin{enumerate}[label=(\Alph*)]
    \item $4$
    \item $6$
    \item $8$
    \item $10$
    \item $12$
\end{enumerate}
\end{exercise}
\begin{solution}
(C). Units of $\Z/20\Z$ are the integers coprime to $20$. Then, use Euler's $\phi$-function.
\end{solution}

\begin{exercise}
Suppose $X$ is compact, $Y$ is Hausdorff, and $f:X\to Y$ is a continuous bijection. Which must be true?
\begin{enumerate}[label=(\Alph*)]
    \item $f^{-1}$ need not be continuous.
    \item $f$ is a homeomorphism.
    \item $Y$ must be compact but $f^{-1}$ can be discontinuous.
    \item $X$ must be connected.
    \item $Y$ must be metric.
\end{enumerate}
\end{exercise}
\begin{solution}
(B). This is a result; a continuous bijection from a compact to Hausdorff has a continuous inverse.
The proof isn't bad, so we'll do this here. Note that the image is compact by continuity,
and surjectivity implies $Y$ is compact. Now, $f$ is a closed map because if you take
some closed $Z \subseteq X$, then it is compact (closed subset of compact is compact), so its image 
is compact, so closed in $Y$. Now, a closed bijeciton is equivalent to an open bijection,
so the inverse is continuous.
\end{solution}

\begin{exercise}
Let $G$ have order $12$. The number $n_3$ of Sylow $3$-subgroups can be
\begin{enumerate}[label=(\Alph*)]
    \item only $1$
    \item only $4$
    \item $1$ or $3$
    \item $1$ or $4$
    \item $3$ or $4$
\end{enumerate}
\end{exercise}
\begin{solution}
(D). Use Sylow III.
\end{solution}

\begin{exercise}
What is the Euler characteristic of the torus $\T^2$?
\begin{enumerate}[label=(\Alph*)]
    \item $-2$
    \item $-1$
    \item $0$
    \item $1$
    \item $2$
\end{enumerate}
\end{exercise}
\begin{solution}
(A). The genus of the $2$-torus (usual torus with one hole) is $1$. This is $0$.
\end{solution}

\begin{remark}
Note that the Euler characteristic is multiplicative across products.
\end{remark}

\begin{exercise}
Over the field $\mathbb{F}_p$, how many roots does $x^p-x$ have?

\begin{enumerate}[label=(\Alph*)]
    \item $1$
    \item $p-1$
    \item $p$
    \item $p+1$
    \item It depends on $p$
\end{enumerate}
\end{exercise}
\begin{solution}
(C). Use FLT.
\end{solution}

\begin{exercise}
The cylinder $\S^1\times[0,1]$ has fundamental group isomorphic to

\begin{enumerate}[label=(\Alph*)]
    \item $0$
    \item $\mathbb{Z}$
    \item $\mathbb{Z}^2$
    \item $\mathbb{Z}/2\mathbb{Z}$
    \item $F_2$
\end{enumerate}
\end{exercise}
\begin{solution}
(B). $[0,1]$ is contractible, so has trivial fundamental group.
$\S^1$ has fundamental group $\Z$.
\end{solution}

\begin{exercise}
Which quotient ring is isomorphic to $\mathbb{C}$?

\begin{enumerate}[label=(\Alph*)]
    \item $\mathbb{R}[x]/(x)$
    \item $\mathbb{R}[x]/(x^2-1)$
    \item $\mathbb{R}[x]/(x^2+1)$
    \item $\mathbb{Z}[x]/(x^2+1)$
    \item $\mathbb{C}[x]/(x^2+1)$
\end{enumerate}
\end{exercise}
\begin{solution}
(C).
\end{solution}

\begin{exercise}
What is the closure of $\mathbb{Q}$ in $\mathbb{R}$ with its usual topology?

\begin{enumerate}[label=(\Alph*)]
    \item $\mathbb{Q}$
    \item $\mathbb{R}\setminus\mathbb{Q}$
    \item $\mathbb{R}$
    \item $\varnothing$
    \item $\mathbb{Q}\cup\{0\}$
\end{enumerate}
\end{exercise}
\begin{solution}
(C). $\Q$ is dense in $\R$ with respect to the usual Euclidean topology 
$d(x,y) = |x - y|$.
\end{solution}

\begin{exercise}
Which of the following spaces is contractible?
\begin{enumerate}[label=(\Alph*)]
    \item $S^1$
    \item $S^2$
    \item $T^2$
    \item $D^2$
    \item $\mathbb{R}^2\setminus\{0\}$
\end{enumerate}
\end{exercise}
\begin{solution}
(D).
\end{solution}

\begin{exercise}
If $G$ acts transitively on a set $X$, and for $x\in X$,
\[
|G|=24,
\qquad
|\operatorname{Stab}(x)|=6,
\]
then $|X|$ is
\begin{enumerate}[label=(\Alph*)]
    \item $3$
    \item $4$
    \item $6$
    \item $18$
    \item $144$
\end{enumerate}
\end{exercise}
\begin{solution}
(B). We use the orbit-stabilizer theorem. The orbit of $x \in X$ has 
cardinality $4$, but transitivity implies that there is only one orbit.
Thus, $|X| = 4$.
\end{solution}

\begin{exercise}
Which statement is true?

\begin{enumerate}[label=(\Alph*)]
    \item Every connected space is path connected.
    \item Every compact space is Hausdorff.
    \item Every path-connected space is connected.
    \item Every bounded subset of a metric space is compact.
    \item Every continuous bijection is a homeomorphism.
\end{enumerate}
\end{exercise}
\begin{solution}
    
\end{solution}

\begin{exercise}
Let
\[
f(z)=\frac{e^z}{z^2(z-1)}.
\]
What is the order of the pole at $z=0$?

\begin{enumerate}[label=(\Alph*)]
    \item Removable
    \item $1$
    \item $2$
    \item $3$
    \item Essential
\end{enumerate}
\end{exercise}

\begin{exercise}
Newton's method for approximating a root of $f(x)=0$ is

\begin{enumerate}[label=(\Alph*)]
    \item $x_{n+1}=x_n-f'(x_n)$
    \item $x_{n+1}=\dfrac{f(x_n)}{f'(x_n)}$
    \item $x_{n+1}=x_n-\dfrac{f(x_n)}{f'(x_n)}$
    \item $x_{n+1}=x_n+\dfrac{f'(x_n)}{f(x_n)}$
    \item $x_{n+1}=f(x_n)$
\end{enumerate}
\end{exercise}

\begin{exercise}
Let $G$ be a finite group and let $H\leq G$ have index $2$. Which must be true?

\begin{enumerate}[label=(\Alph*)]
    \item $H$ is cyclic.
    \item $H$ is abelian.
    \item $H$ is normal.
    \item $G/H$ has four elements.
    \item $G$ is abelian.
\end{enumerate}
\end{exercise}

\begin{exercise}
Which of the following spaces is simply connected?

\begin{enumerate}[label=(\Alph*)]
    \item $S^1$
    \item $\mathbb{R}^2\setminus\{0\}$
    \item $S^2$
    \item $T^2$
    \item An annulus
\end{enumerate}
\end{exercise}

\begin{exercise}
A tree has $17$ vertices. How many edges does it have?

\begin{enumerate}[label=(\Alph*)]
    \item $15$
    \item $16$
    \item $17$
    \item $18$
    \item Cannot be determined
\end{enumerate}
\end{exercise}

\begin{exercise}
Suppose $R$ is an integral domain. Which condition is guaranteed?

\begin{enumerate}[label=(\Alph*)]
    \item Every nonzero element is invertible.
    \item $R$ has no nonzero zero divisors.
    \item Every ideal is principal.
    \item $R$ is finite.
    \item $R$ is a field.
\end{enumerate}
\end{exercise}

\begin{exercise}
Let $\varphi:G\to H$ be a group homomorphism. Which is always a normal subgroup of $G$?

\begin{enumerate}[label=(\Alph*)]
    \item $\operatorname{im}\varphi$
    \item $\ker\varphi$
    \item Every subgroup of $\operatorname{im}\varphi$
    \item Every cyclic subgroup
    \item The center of $H$
\end{enumerate}
\end{exercise}

\begin{exercise}
How many homomorphisms
\[
\mathbb{Z}/6\mathbb{Z}\longrightarrow\mathbb{Z}/4\mathbb{Z}
\]
are there?

\begin{enumerate}[label=(\Alph*)]
    \item $1$
    \item $2$
    \item $3$
    \item $4$
    \item $6$
\end{enumerate}
\end{exercise}

\begin{exercise}
Suppose $f:S^1\to S^1$ is given by $f(z)=z^5$. Under the usual identification $\pi_1(S^1)\cong\mathbb{Z}$, the induced map
\[
f_*:\pi_1(S^1)\to\pi_1(S^1)
\]
is

\begin{enumerate}[label=(\Alph*)]
    \item $n\mapsto 0$
    \item $n\mapsto n$
    \item $n\mapsto 5n$
    \item $n\mapsto n+5$
    \item $n\mapsto n^5$
\end{enumerate}
\end{exercise}

\begin{exercise}
If $F$ is a field, which ideals does $F$ have?

\begin{enumerate}[label=(\Alph*)]
    \item Only $(0)$
    \item Only $F$
    \item Exactly $(0)$ and $F$
    \item Every subset is an ideal
    \item Infinitely many proper ideals
\end{enumerate}
\end{exercise}

\begin{exercise}
Which group is necessarily cyclic?

\begin{enumerate}[label=(\Alph*)]
    \item Every group of order $4$
    \item Every group of order $6$
    \item Every group of prime order
    \item Every abelian group
    \item Every finite group
\end{enumerate}
\end{exercise}